\section*{THE THEORY OF BROWNIAN MOVEMENT.}
\phantomsection
\addcontentsline{toc}{section}{THE THEORY OF BROWNIAN MOVEMENT.}

This theory occupies an exceptional place in contemporary scientific development because of the constant and fruitful exchange which it witnesses between physical problems and ``pure'' mathematics. The study of brownian movement, discovered in 1829 by the botanist Brown, has been carried on intensively in the XIXth century by numerous physicists, \(^{3}\) but the first satisfactory mathematical model was invented by Einstein only in 1905. In the simple case of a particle moving along a straight line, the fundamental hypotheses of Einstein are formulated thus: if \(x(t)\) is the abscissa of the particle at time \(t\) , and if \(t_0 < t_1 < \ldots < t_{n-1} < t_n\) , the successive displacements \(x(t_i) - x(t_{i-1})\) (for \(1 \leq i \leq n\) ) are independent gaussian random variables. It is not the place to evoke here in detail the important experimental work of J. Perrin which motivated Einstein's theory: for our purposes, it is sufficient to remember only a remark of Perrin, according to which the observation of the trajectories of brownian movement suggested irresistibly to him ``the functions without derivative of the mathematicians''. This remark was the initial spark for Wiener.

Quite another stream of ideas draws its origins from the kinetic theory of gases, developed between 1870 and 1900 by Boltzmann and Gibbs. Let us consider a gas made up of N molecules of mass m at (absolute) temperature T and let us denote by \(v_{1}, \ldots, v_{N}\) the velocities of the N molecules of the gas; the kinetic energy of the system is equal to

\[
\frac {m}{2} (\mathbf {v} _ {1} ^ {2} + \dots + \mathbf {v} _ {N} ^ {2}) = 3 N k T\tag{24.1}
\]

where k is Boltzmann's constant. According to the ideas of Gibbs, the multitude of collisions between molecules does not allow the precise determination of the velocities of the molecules, and it is convenient to introduce a probability law P on the sphere S of the space with 3N dimensions defined by the equation (24.1). The ``microcanonical'' hypothesis consists in assuming that P is the measure of mass 1 invariant under rotation on the sphere S. Besides, Maxwell's law of velocities states that the law of probability of the component of the velocity of a molecule has a Gaussian measure of variance 2kT/m. E. Borel seems to have been the first to remark in 1914 that Maxwell's law is a consequence of the hypotheses of Gibbs and of properties of the sphere when the number of molecules is very large. He considers a sphere S in Euclidean space of large dimension and the measure P of mass 1 invariant under rotation on S; using the classical methods of approximation based on the Stirling formula, he shows that the projection of P on a co-ordinate axis is approximately Gaussian. These results are made precise a little later by Gâteaux and Lévy {[}201 a{]}. Given an integer \(m \geq 1\) and a number r > 0, let us denote by \(S_{m,r}\) the set of sequences of the form \((x_{1}, \ldots, x_{m}, 0, 0, \ldots)\) with \(x_{1}^{2} + \ldots + x_{m}^{2} = r^{2}\) ; let us denote also by \(\sigma_{m,r}\) the measure of mass 1 invariant under rotation on \(S_{m,r}\) . Stated in modern language, the result of Gâteaux and Lévy is the following: the sequence of measures \(\sigma_{m,1}\) tends strictly to the unit mass at the origin \((0, 0, \ldots)\) and the sequence of measures \(\sigma_{m,\sqrt{m}}\) tends strictly to a measure \(\Gamma\) of the form

\[
d \Gamma (x _ {1}, x _ {2}, \dots) = \prod_ {n = 1} ^ {\infty} d \gamma (x _ {n})
\]

(γ is the Gaussian measure with variance 1 on R).

The preceding measure \(\Gamma\) plays the role of a Gaussian measure in infinite dimensions. It appears strongly that Lévy had confusedly hoped to define in an intrinsic manner a Gaussian measure on all Hilbert spaces of infinite dimension. In fact, as was shown by Lévy and Wiener, the measure \(\Gamma\) is invariant in a certain sense \(^{4}\) under the automorphisms of \(l^{2}\) ; unfortunately, the set \(l^{2}\) of square summable sequences \((x_{1}, x_{2}, \ldots, x_{n}, \ldots)\) is of measure zero for \(\Gamma\) . It is known today that we must be content with a Gaussian premeasure on a Hilbert space of infinite dimension. \(^{5}\)

To Wiener is owed the essential progress: if there is no reasonable Gaussian measure on a Hilbert space of infinite dimension, it is possible to construct by means of the primitive operation a measure w on a space of continuous functions starting with a Gaussian premeasure. We shall explain succinctly the initial construction of w by Wiener {[}336{]}; it is directly influenced by the relation \(\Gamma = \lim_{m \to \infty} \sigma_{m, \sqrt{m}}\) of Gâteaux and Lévy. For every integer \(m \geq 1\) , let us denote by \(H_m\) the set of functions on T = ]0,1] which are constant on each of the intervals \(\left[\frac{k-1}{m}, \frac{k}{m}\right]\) (for \(k = 1, 2, \ldots, m\) ), and \(\pi_m\) the measure of mass 1 invariant under rotation on the Euclidean sphere of radius 1 in \(R^m\) . Let \(f_m\) be the isomorphism of \(H_m\) on \(R^m\) which associates with each function taking the value \(a_k\) on the interval \(\left[\frac{k-1}{m}, \frac{k}{m}\right]\) the vector \((a_1, a_2 - a_1, \ldots, a_m - a_{m-1})\) (whence the name of ``differential space'' dear to Wiener); let us denote by \(w_m\) the measure on \(H_m\) image of \(\pi_m\) under \(f_m^{-1}\) . Wiener defines the measure w being sought as the limit of the measures \(w_m\) . Precisely, let us denote by H the set of regular functions on T, with the uniform convergence topology (we have \(H_m \subset H\) for every integer \(m \geq 1\) ); for every uniformly continuous function F bounded on H, the limit \(A\{F\} = \lim_{m \to \infty} \int_{H_m} F(x) dw_m(x)\) exists; Wiener then obtains certain majorations by a subtle analysis of the fluctuations of the items in play, and, taking up again the arguments of compactness brought out by Daniell, he shows that we are in a situation for the application of Daniell's extension theorem. We conclude the existence of a measure w carried by \(C(T)\) and such that \(A\{F\} = \int_{C(T)} F(x) dw(x)\) . Wiener can then show that the measure w corresponds to the hypotheses of Einstein, \(^{6}\) and his estimates allowed him to give a precise meaning to the remark of Perrin on functions without a derivative: the set of functions satisfying a Lipschitz condition of order \(\frac{1}{2}\) , is negligible for w (on the contrary, for every a with \(0 < a < \frac{1}{2}\) , nearly every function satisfies a Lipschitz condition of order a).

We know today numerous constructions for the Wiener measure. Thus, Paley and Wiener use random Fourier series ({[}243{]}, chap.~IX): to every real sequence \(\mathbf{a} = (a_{n})_{n \geq 1}\) and every integer \(m \geq 0\) , let us define the function \(f_{m,a}\) on ]0, 1] by

\[
f _ {m, \mathbf {a}} (t) = a _ {1} t + 2 \sum_ {k = 2} ^ {2 ^ {m + 1}} \frac {1}{\pi k} a _ {k - 1} \sin \pi k t;
\]

it can be shown that, for \(\Gamma\) -almost all sequences \(\mathbf{a}\) , the sequence of functions \(f_{m,\mathbf{a}}\) tends to a continuous function \(f_{\mathbf{a}}\) and that \(w\) is the image of \(\Gamma\) by the map (defined almost everywhere) \(\mathbf{a} \mapsto f_{\mathbf{a}}\) . Later, Lévy gave in {[}201 b and c{]} another construction. Finally Kac, Donsker and Erdős show around 1950 how to replace in Wiener's initial construction the spherical measures \(\pi_m\) on \(\mathbf{R}^m\) by more general measures. Their results establish a solid link between Wiener's measure and the limit theorems of Probability Theory; they will be completed and systematised by Prokhorov {[}254{]} in work to which we will return later.

This is not the place to analyse the numerous and important probabilistic works caused by Wiener's discovery; today, the brownian movement appears only as one of the most important examples of a Markov process. We will only mention the application made by Kac of Wiener's measure to the solution of certain parabolic partial differential equations; it is a matter there of an adaptation of the ideas of Feynman in quantum field theory --- another example of this reciprocal influence between mathematics and the problems of physics.

